\subsection{Zielsetzung und Verengung}\label{sec:zielsetzung_und_verengung}

Aus der beschriebenen Lücke ergibt sich das Ziel dieses Projektberichts: die Konzeption der Plattform \enquote{Resteria}. Resteria soll Reste-Matching, also Rezeptvorschläge auf Basis vorhandener Zutaten, mit einer echten Social-Community verbinden, die auf Zero-Waste-Kochen ausgerichtet ist. Nutzer:innen sollen dort nicht nur passiv Rezepte finden, sondern sich über ihre Reste-Kreationen austauschen, sich gegenseitig inspirieren und gemeinsam an Herausforderungen wie einer Zero-Waste-Woche teilnehmen.

Dieses Konzept verengt die Zielgruppe bewusst auf nachhaltigkeitsorientierte Hobbyköch:innen, also Menschen, die beim Kochen bereits Wert auf Resteverwertung, saisonale oder regionale Zutaten legen oder sich dafür interessieren. Diese Verengung ist nötig, weil eine Plattform, die alle Hobbyköch:innen gleichermaßen ansprechen will, ihre Kernfunktionen zwangsläufig verwässert. Gamification-Elemente wie Rettungspunkte oder Reste-Level ergeben nur für Menschen einen Anreiz, denen Nachhaltigkeit beim Kochen bereits wichtig ist. Für die breite Masse der Hobbyköch:innen, die vor allem neue Gerichte entdecken wollen, wäre ein solcher Fokus ohne echten Mehrwert.

Diese Zielgruppenannahme ist argumentativ begründet, nicht empirisch mit eigenen Nutzerdaten geprüft.
